\documentclass[10pt]{article}

\usepackage[margin=0.9in]{geometry}
\usepackage{enumitem}
\usepackage{booktabs}

\setlength{\parskip}{0.25em}
\setlist{nosep, leftmargin=*}

\title{\vspace{-1.5em}ECE 268 Project Progress Report\\[0.2em]
\large Lightweight Block Ciphers for Constrained Devices}
\author{%
  Narain Shriraam M S (A69041741) \and
  Karthik N R \and
  Tushar Nasery
}
\date{June 3, 2026}

\begin{document}
\maketitle
\vspace{-1.2em}

\section*{Project Summary}
\textbf{Goal:} Implement SIMON-64/128 and GIFT-128/128 with encryption, decryption, key schedules, and CBC/CTR modes; validate against published test vectors; benchmark against an AES-128 baseline (code size, cycles/byte, key-schedule cost) for IoT/RFID settings where AES is too costly. \\
\textbf{Status:} All implementations are complete and correct. 16/16 test cases pass including spec-mandated ECB vectors. Benchmarks and code-size measurements are collected.

\section{Work Completed}

\subsection*{Implementations}
\begin{itemize}
  \item \textbf{Repository layout:} \texttt{simon64.c/h}, \texttt{gift128.c/h}, \texttt{modes.c/h}, \texttt{test\_vectors.c}, \texttt{bench.c}, and \texttt{Makefile}.
  \item \textbf{SIMON 64/128:} 64-bit block, 128-bit key, 44 rounds (ePrint 2013/404). Uses $z_3$ key schedule (m=4), big-endian block I/O. Feistel round function uses circular shifts (1, 2, 8) and bitwise AND/XOR.
  \item \textbf{GIFT-128/128:} 128-bit block/key, 40 rounds (ePrint 2017/622). Nibble-based SPN with 4-bit S-box, 128-bit bit-permutation, 6-bit round constants, and key update. Precomputed round keys \texttt{rk[40][2]}.
  \item \textbf{Modes:} Cipher-agnostic CBC (PKCS\#7 padding) and CTR (big-endian counter) via function-pointer interface \texttt{block\_fn\_t}. Any cipher plugs in without modification.
  \item \textbf{Test harness:} 16 test cases covering SIMON ECB, GIFT ECB (2 spec vectors), CBC round-trips for both ciphers, CTR round-trips, and PKCS\#7 padding edge cases (1-byte, block-aligned).
  \item \textbf{Benchmark:} \texttt{bench.c} measures key-schedule cost (ns), ECB and CTR throughput (cycles/byte) for SIMON, GIFT, and AES-128 (CommonCrypto hw-accelerated baseline).
\end{itemize}

\subsection*{Test Vector Validation}
\begin{table}[h]
\centering\small
\begin{tabular}{lcccc}
\toprule
\textbf{Cipher} & \textbf{Block / Key} & \textbf{Spec vectors} & \textbf{CBC/CTR} & \textbf{PKCS\#7} \\
\midrule
SIMON 64/128 & 64 / 128 & \textbf{Pass} & \textbf{Pass} & \textbf{Pass} \\
GIFT-128/128 & 128 / 128 & \textbf{Pass (2/2)} & \textbf{Pass} & --- \\
\bottomrule
\end{tabular}
\end{table}

\subsection*{Benchmark Results (Apple Silicon, \texttt{-O2}, 64 KiB)}
\begin{table}[h]
\centering\small
\begin{tabular}{lcccr}
\toprule
\textbf{Cipher} & \textbf{Key Sched (ns)} & \textbf{ECB (cpb)} & \textbf{CTR (cpb)} & \textbf{.text (B)} \\
\midrule
SIMON 64/128 & 38 & 10.9 & 10.8 & 356 \\
GIFT-128/128 & 488 & 358.8 & 338.9 & 4\,476 \\
AES-128 (hw) & 448 & 1.8 & 1.6 & --- \\
\bottomrule
\end{tabular}
\end{table}
\textbf{Analysis:} SIMON's Feistel structure with simple bitwise operations yields a tiny code footprint (356 B) and competitive throughput ($\approx$11 cpb). GIFT's nibble-based SPN permutation is significantly slower in software ($\approx$340 cpb) but is designed for minimal hardware gate count ($\sim$1300 GE per literature). AES-128 with hardware AES-NI achieves $<$2 cpb but is unavailable on the target constrained platforms.

\section{Challenges and Issues Encountered}
\begin{itemize}
  \item \textbf{SIMON variant confusion:} The original implementation used SIMON 64/64 (m=2, 32 rounds, $z_0$), a non-standard variant absent from the spec appendix. Corrected to SIMON 64/128 (m=4, 44 rounds, $z_3$) which has a published test vector.
  \item \textbf{GIFT nibble packing:} The byte-to-nibble conversion had swapped high/low nibbles, causing all encrypt outputs to mismatch. Fixed by aligning with the reference C++ implementation's ordering.
  \item \textbf{Key schedule constants:} Ensuring the correct $z$-sequence constant (\texttt{0xfc2ce51207a635db}) for $z_3$ required cross-checking against Crypto++ and refinery implementations.
\end{itemize}

\section{Plans for Final Submission}
\textbf{Remaining work:}
\begin{enumerate}
  \item Final report: expand analysis with gate-count references, Feistel vs.\ SPN trade-off discussion.
  \item Presentation: 10--12 slides with architecture diagrams and benchmark bar charts.
  \item Stretch: compare at additional data sizes (1 KB, 16 KB) and add a third cipher (SPECK-64 or PRESENT-80).
\end{enumerate}

\begin{table}[h]
\centering\small
\begin{tabular}{p{0.22\linewidth}p{0.7\linewidth}}
\toprule
\textbf{Member} & \textbf{Tasks} \\
\midrule
Narain Shriraam M S & GIFT-128 implementation, benchmark scripts, final report \\
Karthik N R & SIMON-64 implementation, AES baseline, results plots \\
Tushar Nasery & Test harness, Makefile, modes, slides \\
\bottomrule
\end{tabular}
\end{table}

\section*{Feasibility and Direction}
The project is on track. All core deliverables (two cipher implementations, CBC/CTR wrappers, test vector validation, and AES-baseline benchmark) are complete and passing. The remaining work is report writing and presentation preparation. The benchmark data confirms the design rationale: SIMON offers the best software performance among lightweight ciphers, while GIFT targets hardware efficiency.

\vspace{0.3em}
\small\textit{References:} Bogdanov et al., PRESENT (CHES 2007); Beaulieu et al., SIMON/SPECK (ePrint 2013/404); Banik et al., GIFT (ePrint 2017/622).

\end{document}
